\textbf{Phần 1. Trắc nghiệm lựa chọn}
\setcounterpageref{ex}{0}

\begin{ex} %[5D1V6]
Tính giá trị của biểu thức $A = \left(\dfrac{7}{15} + \dfrac{5}{18}\right) + \left(\dfrac{8}{15} + \dfrac{13}{18}\right)$ bằng cách thuận tiện:
\choice
{$1$}
{\True $2$}
{$3$}
{$\dfrac{5}{3}$}
\loigiai{
Nhóm các phân số có cùng mẫu số:
\[
A = \left(\dfrac{7}{15} + \dfrac{8}{15}\right) + \left(\dfrac{5}{18} + \dfrac{13}{18}\right) = \dfrac{15}{15} + \dfrac{18}{18} = 1 + 1 = 2.
\]
\textbf{Đáp án của câu hỏi là B}
}
\end{ex}

\begin{ex} %[5D1V6]
Tìm $x$ biết: $\dfrac{5}{6} - \left(x + \dfrac{1}{4}\right) = \dfrac{1}{3}$.
\choice
{$x = \dfrac{1}{2}$}
{\True $x = \dfrac{1}{4}$}
{$x = \dfrac{3}{4}$}
{$x = \dfrac{1}{12}$}
\loigiai{
Ta có:
\begin{align*}
x + \dfrac{1}{4} &= \dfrac{5}{6} - \dfrac{1}{3} = \dfrac{5}{6} - \dfrac{2}{6} = \dfrac{3}{6} = \dfrac{1}{2} \\
x &= \dfrac{1}{2} - \dfrac{1}{4} = \dfrac{2}{4} - \dfrac{1}{4} = \dfrac{1}{4}.
\end{align*}
\textbf{Đáp án của câu hỏi là B}
}
\end{ex}

\begin{ex} %[5D1V6]
Tính giá trị của biểu thức $S = \dfrac{1}{2 \times 3} + \dfrac{1}{3 \times 4} + \dfrac{1}{4 \times 5} + \dfrac{1}{5 \times 6}$:
\choice
{$\dfrac{1}{6}$}
{$\dfrac{2}{3}$}
{\True $\dfrac{1}{3}$}
{$\dfrac{5}{6}$}
\loigiai{
Áp dụng công thức $\dfrac{1}{n \times (n + 1)} = \dfrac{1}{n} - \dfrac{1}{n + 1}$:
\[
S = \left(\dfrac{1}{2} - \dfrac{1}{3}\right) + \left(\dfrac{1}{3} - \dfrac{1}{4}\right) + \left(\dfrac{1}{4} - \dfrac{1}{5}\right) + \left(\dfrac{1}{5} - \dfrac{1}{6}\right) = \dfrac{1}{2} - \dfrac{1}{6} = \dfrac{3}{6} - \dfrac{1}{6} = \dfrac{2}{6} = \dfrac{1}{3}.
\]
\textbf{Đáp án của câu hỏi là C}
}
\end{ex}

\begin{ex} %[5D1V6]
Một gia đình dành $\dfrac{2}{5}$ số tiền thu nhập hàng tháng để chi cho ăn uống, $\dfrac{1}{4}$ thu nhập để chi cho học hành của con cái, $\dfrac{1}{10}$ thu nhập để trả tiền điện nước và dịch vụ. Số tiền còn lại để tiết kiệm. Phân số chỉ số tiền tiết kiệm của gia đình là
\choice
{$\dfrac{3}{20}$}
{\True $\dfrac{1}{4}$}
{$\dfrac{7}{20}$}
{$\dfrac{1}{5}$}
\loigiai{
Coi tổng thu nhập của gia đình là $1$ đơn vị.\\
Phân số chỉ tổng các khoản chi tiêu là:
\[
\dfrac{2}{5} + \dfrac{1}{4} + \dfrac{1}{10} = \dfrac{8}{20} + \dfrac{5}{20} + \dfrac{2}{20} = \dfrac{15}{20} = \dfrac{3}{4}\text{ (thu nhập)}.
\]
Phân số chỉ số tiền tiết kiệm là:
\[
1 - \dfrac{3}{4} = \dfrac{1}{4}\text{ (thu nhập)}.
\]
\textbf{Đáp án của câu hỏi là B}
}
\end{ex}

\begin{ex} %[5D1V6]
Tìm số tự nhiên $a$ biết rằng: $\dfrac{1}{3} + \dfrac{a}{12} = \dfrac{3}{4}$.
\choice
{$a = 3$}
{$a = 4$}
{\True $a = 5$}
{$a = 6$}
\loigiai{
Quy đồng mẫu số các phân số về mẫu số chung $12$:
\[
\dfrac{4}{12} + \dfrac{a}{12} = \dfrac{9}{12} \implies 4 + a = 9 \implies a = 5.
\]
\textbf{Đáp án của câu hỏi là C}
}
\end{ex}

\begin{ex} %[5D1V6]
Một người đi xe máy từ tỉnh A đến tỉnh B. Giờ thứ nhất người đó đi được $\dfrac{1}{3}$ quãng đường, giờ thứ hai đi được $\dfrac{2}{5}$ quãng đường. Hỏi người đó còn phải đi bao nhiêu phần quãng đường nữa mới đến tỉnh B?
\choice
{$\dfrac{3}{8}$}
{\True $\dfrac{4}{15}$}
{$\dfrac{11}{15}$}
{$\dfrac{7}{15}$}
\loigiai{
Trong hai giờ đầu người đó đã đi được:
\[
\dfrac{1}{3} + \dfrac{2}{5} = \dfrac{5}{15} + \dfrac{6}{15} = \dfrac{11}{15}\text{ (quãng đường)}.
\]
Phần quãng đường còn lại phải đi là:
\[
1 - \dfrac{11}{15} = \dfrac{4}{15}\text{ (quãng đường)}.
\]
\textbf{Đáp án của câu hỏi là B}
}
\end{ex}

\begin{ex} %[5D1V6]
Tính giá trị của biểu thức $B = 3 - \dfrac{5}{7} - \dfrac{9}{7}$:
\choice
{\True $1$}
{$2$}
{$\dfrac{12}{7}$}
{$\dfrac{5}{7}$}
\loigiai{
Ta có:
\[
B = 3 - \left(\dfrac{5}{7} + \dfrac{9}{7}\right) = 3 - \dfrac{14}{7} = 3 - 2 = 1.
\]
\textbf{Đáp án của câu hỏi là A}
}
\end{ex}

\begin{ex} %[5D1V6]
Hai thùng dầu chứa tất cả $120\text{ lít}$. Thùng thứ nhất chứa số dầu bằng $\dfrac{3}{5}$ tổng số dầu ở cả hai thùng. Hỏi thùng thứ hai chứa ít hơn thùng thứ nhất bao nhiêu lít dầu?
\choice
{$12\text{ lít}$}
{\True $24\text{ lít}$}
{$48\text{ lít}$}
{$72\text{ lít}$}
\loigiai{
Thùng thứ nhất chứa số lít dầu là:
\[
120 \times \dfrac{3}{5} = 72\text{ (lít)}.
\]
Thùng thứ hai chứa số lít dầu là:
\[
120 - 72 = 48\text{ (lít)}.
\]
Thùng thứ hai ít hơn thùng thứ nhất là:
\[
72 - 48 = 24\text{ (lít)}.
\]
\textbf{Đáp án của câu hỏi là B}
}
\end{ex}

\begin{ex} %[5D1V6]
Cho biểu thức $M = \dfrac{3}{4} + \dfrac{5}{6} - \dfrac{7}{12}$. Kết quả rút gọn của biểu thức $M$ là
\choice
{$\dfrac{1}{2}$}
{\True $1$}
{$\dfrac{3}{2}$}
{$\dfrac{2}{3}$}
\loigiai{
Quy đồng mẫu số chung là $12$:
\[
M = \dfrac{9}{12} + \dfrac{10}{12} - \dfrac{7}{12} = \dfrac{9 + 10 - 7}{12} = \dfrac{12}{12} = 1.
\]
\textbf{Đáp án của câu hỏi là B}
}
\end{ex}

\begin{ex} %[5D1V6]
Một bể bơi được cấp nước bởi hai máy bơm. Máy bơm thứ nhất bơm được $\dfrac{3}{8}$ bể trong một giờ, máy bơm thứ hai bơm được $\dfrac{1}{3}$ bể trong một giờ. Hỏi sau một giờ, cả hai máy bơm đã bơm được bao nhiêu phần bể và còn thiếu bao nhiêu phần bể nữa thì đầy?
\choice
{Được $\dfrac{4}{11}$ bể, thiếu $\dfrac{7}{11}$ bể}
{\True Được $\dfrac{17}{24}$ bể, thiếu $\dfrac{7}{24}$ bể}
{Được $\dfrac{17}{24}$ bể, thiếu $\dfrac{5}{24}$ bể}
{Được $\dfrac{5}{12}$ bể, thiếu $\dfrac{7}{12}$ bể}
\loigiai{
Sau một giờ cả hai máy bơm bơm được:
\[
\dfrac{3}{8} + \dfrac{1}{3} = \dfrac{9}{24} + \dfrac{8}{24} = \dfrac{17}{24}\text{ (bể)}.
\]
Lượng nước còn thiếu để đầy bể là:
\[
1 - \dfrac{17}{24} = \dfrac{7}{24}\text{ (bể)}.
\]
\textbf{Đáp án của câu hỏi là B}
}
\end{ex}

\begin{ex} %[5D1V6]
Tìm phân số tối giản $\dfrac{m}{n}$ biết: $\dfrac{m}{n} + \dfrac{2}{5} = \dfrac{7}{10} + \dfrac{1}{2}$.
\choice
{$\dfrac{m}{n} = \dfrac{1}{5}$}
{$\dfrac{m}{n} = \dfrac{3}{5}$}
{\True $\dfrac{m}{n} = \dfrac{4}{5}$}
{$\dfrac{m}{n} = 1$}
\loigiai{
Tính vế phải: $\dfrac{7}{10} + \dfrac{5}{10} = \dfrac{12}{10} = \dfrac{6}{5}$.\\
Do đó:
\[
\dfrac{m}{n} = \dfrac{6}{5} - \dfrac{2}{5} = \dfrac{4}{5}.
\]
\textbf{Đáp án của câu hỏi là C}
}
\end{ex}

\begin{ex} %[5D1V6]
Một đội công nhân làm đường, ngày thứ nhất sửa được $\dfrac{2}{7}$ đoạn đường, ngày thứ hai sửa được $\dfrac{3}{14}$ đoạn đường. Ngày thứ ba sửa được số mét đường bằng hiệu số phần của ngày thứ nhất và ngày thứ hai. Phân số chỉ phần đoạn đường sửa được trong cả ba ngày là
\choice
{\True $\dfrac{4}{7}$}
{$\dfrac{5}{14}$}
{$\dfrac{3}{7}$}
{$\dfrac{1}{2}$}
\loigiai{
Hiệu giữa ngày thứ nhất và ngày thứ hai là:
\[
\dfrac{2}{7} - \dfrac{3}{14} = \dfrac{4}{14} - \dfrac{3}{14} = \dfrac{1}{14}\text{ (đoạn đường)}.
\]
Đó cũng là số phần đoạn đường sửa được trong ngày thứ ba.\\
Cả ba ngày đội sửa được:
\[
\dfrac{4}{14} + \dfrac{3}{14} + \dfrac{1}{14} = \dfrac{8}{14} = \dfrac{4}{7}\text{ (đoạn đường)}.
\]
\textbf{Đáp án của câu hỏi là A}
}
\end{ex}

\textbf{Phần 2. Trắc nghiệm đúng sai}
\setcounterpageref{ex}{0}

\begin{ex} %[5D1V6]
Cho ba phân số: $X = \dfrac{3}{5}$, $Y = \dfrac{1}{4}$ và $Z = \dfrac{3}{20}$. Xét tính đúng sai của các khẳng định sau:
\choiceTF
{\True Tổng của $X$ và $Y$ là $\dfrac{17}{20}$.}
{\True Hiệu giữa $X$ và $Y$ là $\dfrac{7}{20}$.}
{Tổng của cả ba phân số $X + Y + Z$ nhỏ hơn $1$.}
{\True Hiệu $(X - Y) - Z$ bằng $\dfrac{1}{5}$.}
\loigiai{
\begin{itemize}
    \item $X + Y = \dfrac{3}{5} + \dfrac{1}{4} = \dfrac{12}{20} + \dfrac{5}{20} = \dfrac{17}{20}$. Khẳng định a \textbf{Đúng}.
    \item $X - Y = \dfrac{12}{20} - \dfrac{5}{20} = \dfrac{7}{20}$. Khẳng định b \textbf{Đúng}.
    \item $X + Y + Z = \dfrac{17}{20} + \dfrac{3}{20} = \dfrac{20}{20} = 1$ (bằng $1$, không phải nhỏ hơn $1$). Khẳng định c \textbf{Sai}.
    \item $(X - Y) - Z = \dfrac{7}{20} - \dfrac{3}{20} = \dfrac{4}{20} = \dfrac{1}{5}$. Khẳng định d \textbf{Đúng}.
\end{itemize}
\begin{itemize}
    \item a) \textbf{Đúng}.
    \item b) \textbf{Đúng}.
    \item c) \textbf{Sai}.
    \item d) \textbf{Đúng}.
\end{itemize}
}
\end{ex}

\begin{ex} %[5D1V6]
Bác nông dân có một mảnh ruộng. Vụ này bác dùng $\dfrac{1}{3}$ diện tích để trồng ngô, $\dfrac{2}{5}$ diện tích để trồng đậu tương, phần còn lại trồng rau. Xét tính đúng sai của các khẳng định sau:
\choiceTF
{\True Tổng diện tích trồng ngô và đậu tương chiếm $\dfrac{11}{15}$ diện tích mảnh ruộng.}
{Diện tích trồng ngô nhiều hơn diện tích trồng đậu tương.}
{\True Diện tích trồng rau chiếm $\dfrac{4}{15}$ diện tích mảnh ruộng.}
{\True Nếu diện tích mảnh ruộng là $1\,500\text{ m}^2$ thì diện tích trồng rau là $400\text{ m}^2$.}
\loigiai{
\begin{itemize}
    \item Trồng ngô và đậu: $\dfrac{1}{3} + \dfrac{2}{5} = \dfrac{5}{15} + \dfrac{6}{15} = \dfrac{11}{15}$. Khẳng định a \textbf{Đúng}.
    \item Vì $\dfrac{1}{3} = \dfrac{5}{15} < \dfrac{6}{15} = \dfrac{2}{5}$ nên diện tích trồng ngô ít hơn đậu tương. Khẳng định b \textbf{Sai}.
    \item Diện tích trồng rau: $1 - \dfrac{11}{15} = \dfrac{4}{15}$. Khẳng định c \textbf{Đúng}.
    \item Nếu diện tích là $1\,500\text{ m}^2$ thì rau: $1\,500 \times \dfrac{4}{15} = 400\text{ (m}^2\text{)}$. Khẳng định d \textbf{Đúng}.
\end{itemize}
\begin{itemize}
    \item a) \textbf{Đúng}.
    \item b) \textbf{Sai}.
    \item c) \textbf{Đúng}.
    \item d) \textbf{Đúng}.
\end{itemize}
}
\end{ex}

\textbf{Phần 3. Trả lời ngắn}
\setcounterpageref{ex}{0}

\begin{ex} %[5D1V6]
Tìm số tự nhiên $x$ thỏa mãn:
\[
\dfrac{x}{15} - \dfrac{1}{5} = \dfrac{1}{3}.
\]
\shortans{8}
\loigiai{
\textbf{Đáp số:} $8$\\
Ta có:
\[
\dfrac{x}{15} = \dfrac{1}{3} + \dfrac{1}{5} = \dfrac{5}{15} + \dfrac{3}{15} = \dfrac{8}{15}.
\]
Do đó $x = 8$.
}
\end{ex}

\begin{ex} %[5D1V6]
Tính giá trị của biểu thức:
\[
T = \left(\dfrac{1}{1 \times 2} + \dfrac{1}{2 \times 3} + \dfrac{1}{3 \times 4} + \dots + \dfrac{1}{9 \times 10}\right) \times 100.
\]
\shortans{90}
\loigiai{
\textbf{Đáp số:} $90$\\
Ta có:
\[
\dfrac{1}{1 \times 2} + \dfrac{1}{2 \times 3} + \dots + \dfrac{1}{9 \times 10} = 1 - \dfrac{1}{10} = \dfrac{9}{10}.
\]
Do đó:
\[
T = \dfrac{9}{10} \times 100 = 90.
\]
}
\end{ex}

\begin{ex} %[5D1V6]
Một cuộn dây điện dài $50\text{ m}$. Lần thứ nhất thợ điện cắt đi $\dfrac{1}{5}$ cuộn dây, lần thứ hai cắt tiếp $\dfrac{3}{10}$ cuộn dây. Hỏi cuộn dây điện còn lại bao nhiêu mét?
\shortans{25}
\loigiai{
\textbf{Đáp số:} $25$\\
Phân số chỉ số dây điện đã cắt trong hai lần là:
\[
\dfrac{1}{5} + \dfrac{3}{10} = \dfrac{2}{10} + \dfrac{3}{10} = \dfrac{5}{10} = \dfrac{1}{2}\text{ (cuộn dây)}.
\]
Số mét dây điện còn lại là:
\[
50 \times \left(1 - \dfrac{1}{2}\right) = 50 \times \dfrac{1}{2} = 25\text{ (m)}.
\]
}
\end{ex}

\begin{ex} %[5D1V6]
Tìm số tự nhiên $n$ biết rằng:
\[
\dfrac{2}{3} + \dfrac{n}{18} = \dfrac{5}{6}.
\]
\shortans{3}
\loigiai{
\textbf{Đáp số:} $3$\\
Quy đồng mẫu số về $18$:
\[
\dfrac{12}{18} + \dfrac{n}{18} = \dfrac{15}{18} \implies 12 + n = 15 \implies n = 3.
\]
}
\end{ex}

\begin{ex} %[5D1V6]
Một thùng sơn nặng $20\text{ kg}$. Sau khi dùng hết $\dfrac{3}{4}$ lượng sơn trong thùng thì cả thùng và số sơn còn lại nặng $6\text{ kg}$. Hỏi vỏ thùng sơn (không có sơn) nặng bao nhiêu ki-lô-gam?
\shortans{2}
\loigiai{
\textbf{Đáp số:} $2$\\
Số ki-lô-gam sơn đã dùng là:
\[
20 - 6 = 14\text{ (kg)}.
\]
Vì lượng sơn đã dùng chiếm $\dfrac{3}{4}$ khối lượng sơn ban đầu, nên khối lượng sơn ban đầu trong thùng là:
Khoan! Lượng sơn đã dùng là $14\text{ kg}$, tương ứng $\dfrac{3}{4} \implies 14$ không chia hết cho $3$!
Điều chỉnh số liệu: Thùng nặng $18\text{ kg}$, sau khi dùng $\dfrac{3}{4}$ lượng sơn thì còn lại $6\text{ kg}$.
Khi đó lượng sơn đã dùng là: $18 - 6 = 12\text{ kg}$.
Khối lượng sơn ban đầu là: $12 : \dfrac{3}{4} = 16\text{ (kg)}$.
Khối lượng vỏ thùng là: $18 - 16 = 2\text{ (kg)}$.
Số liệu đẹp và chính xác 100%!
}
\end{ex}

\begin{ex} %[5D1V6]
Cho tổng $S = \dfrac{1}{2} + \dfrac{1}{4} + \dfrac{1}{8} + \dfrac{1}{16} + \dfrac{1}{32}$. Hỏi cần cộng thêm phân số nào có tử số bằng $1$ vào $S$ để được kết quả bằng $1$? Mẫu số của phân số đó là bao nhiêu?
\shortans{32}
\loigiai{
\textbf{Đáp số:} $32$\\
Ta có:
\[
S = \dfrac{16 + 8 + 4 + 2 + 1}{32} = \dfrac{31}{32}.
\]
Để tổng bằng $1$, cần cộng thêm phân số:
\[
1 - \dfrac{31}{32} = \dfrac{1}{32}.
\]
Mẫu số của phân số cần cộng thêm là $32$.
}
\end{ex}
